% Deutsche Parallelfassung zu foundations/core_equations.tex

\chapter{Kerngleichungen}
\label{chap:core-equations}

Die folgenden Gleichungen sind ausführbares Wissen: Jede Beziehung benennt die
Eingaben, die sie empfängt, die Ausgabe, die sie erzeugt, die Annahmen, unter
denen diese Ausgabe Bedeutung besitzt, und die Fragen, die sie nicht
entscheiden kann. Die Folge begleitet den wiederkehrenden Fall von CAD-Kurve
und Vermessungspolyline von der Konstruktion über Realisierung und
Repräsentation bis zum Vergleich.

\section{Intrinsischer Zustand und Steuerungen}

Der grundlegende räumliche Eisenbahnzustand wird geschrieben als
\begin{equation}
x(s)
=
\bigl(\gamma(s),\theta(s),\phi(s)\bigr),
\label{eq:canonical-state}
\end{equation}
wobei \(\gamma(s)\) die ebene Position, \(\theta(s)\) den Richtungswinkel und
\(\phi(s)\) die Überhöhungsdrehung bezeichnet.

Das zugehörige Steuerungspaar lautet
\begin{equation}
u(s)
=
\bigl(\kappa(s),\omega(s)\bigr),
\qquad
\omega(s)=\phi'(s).
\label{eq:control-variables}
\end{equation}

Die Unterscheidung von Zustand und Steuerung folgt dem Zustandsmodell aus
\cref{chap:state-model}: Der Zustand hält den Zustand des Alignments fest,
während die Steuerungen seine Entwicklung entlang der Bogenlänge bestimmen.

\paragraph{Ausführungsvertrag.}
Die Gleichungen empfangen eine longitudinale Domäne, einen Anfangszustand und
zulässige Steuerungen. Sie erzeugen eine Zustandstrajektorie. Sie setzen
voraus, dass das gewählte Zustandsmodell und die Regularitätsbedingungen
gelten. Sie identifizieren weder das Engineering Object noch bestimmen sie die
Autorität einer Datenquelle.

\section{Ebene intrinsische Geometrie}

Die Bogenlängenparametrisierung ergibt
\begin{equation}
\gamma'(s)=\mathbf t(s),
\label{eq:state-evolution}
\end{equation}
und die Richtungsentwicklung wird bestimmt durch
\begin{equation}
\theta'(s)=\kappa(s).
\label{eq:curvature-relation}
\end{equation}

Im ebenen Fall ist der Einheitstangentenvektor
\begin{equation}
\mathbf t(s)
=
\begin{pmatrix}
\cos\theta(s)\\
\sin\theta(s)
\end{pmatrix}.
\label{eq:planar-tangent}
\end{equation}

Das kanonische ebene Rekonstruktionssystem lautet somit
\begin{equation}
\begin{aligned}
\gamma'(s)
&=
\begin{pmatrix}
\cos\theta(s)\\
\sin\theta(s)
\end{pmatrix},\\
\theta'(s)&=\kappa(s).
\end{aligned}
\label{eq:pose2-system}
\end{equation}

Bei gegebener Anfangsposition und gegebenem Anfangswinkel bestimmt die
Krümmungsfunktion die ebene Alignment-Trajektorie.

\paragraph{Ausführungsvertrag.}
Die Rekonstruktion empfängt \(\kappa(s)\), \(\gamma(s_0)\) und
\(\theta(s_0)\); sie erzeugt \(\theta(s)\), \(\mathbf t(s)\) und
\(\gamma(s)\). Sie setzt Bogenlängenparametrisierung und hinreichende
Regularität voraus. Eine aus \(\gamma\) abgetastete CAD-Kurve ist daher eine
abgeleitete Repräsentation. Die Rekonstruktion kann weder entscheiden, ob eine
Vermessungspolyline dasselbe Alignment beobachtet, noch ob eine Abweichung
akzeptabel ist.

\section{Überhöhungsentwicklung und grundlegendes pose3-System}

Die Überhöhungsdrehung entwickelt sich gemäß
\begin{equation}
\phi'(s)=\omega(s).
\label{eq:cant-evolution}
\end{equation}

Zusammen mit den ebenen Rekonstruktionsgleichungen ergibt sich das
grundlegende pose3-System
\begin{equation}
\begin{aligned}
\gamma'(s)
&=
\begin{pmatrix}
\cos\theta(s)\\
\sin\theta(s)
\end{pmatrix},\\
\theta'(s)&=\kappa(s),\\
\phi'(s)&=\omega(s).
\end{aligned}
\label{eq:pose3-system}
\end{equation}

Abstrakt kann die Zustandsentwicklung geschrieben werden als
\begin{equation}
x'(s)=f\bigl(x(s),u(s)\bigr).
\label{eq:general-state-system}
\end{equation}

Diese Gleichung drückt ausschließlich die Zustands-System-Form aus. Profil,
metrische Realisierung, räumliche Frame-Konstruktion und physikalische
Interpretation werden getrennt eingeführt.

\paragraph{Ausführungsvertrag.}
Das System empfängt Krümmungs- und Überhöhungsratensteuerungen sowie einen
Anfangszustand; es erzeugt eine Pose-Trajektorie. Es setzt voraus, dass das
gewählte Pose-Modell anwendbar ist. Es erzeugt weder eine physische
Realisierung noch deren Beobachtung.

\section{Dynamische Eisenbahnbeziehungen}

Für die klassische lokale Ingenieurnäherung lautet die effektive
Querbeschleunigung
\begin{equation}
a_{\mathrm{eff}}
=
v^2\kappa-g\sin\phi.
\label{eq:effective-lateral-acceleration}
\end{equation}

Idealer Überhöhungsausgleich erfüllt
\begin{equation}
v^2\kappa=g\sin\phi.
\label{eq:cant-equilibrium}
\end{equation}

Bei konstanter Geschwindigkeit ergibt die Differentiation nach der Zeit die
zugehörige Ruckbeziehung
\begin{equation}
j
=
v^3\kappa'
-gv\cos\phi\,\phi'.
\label{eq:jerk-relation}
\end{equation}

Diese Gleichungen liefern die grundlegende Kopplung zwischen Krümmung,
Überhöhung und betrieblicher Geschwindigkeit. Ihre eisenbahntechnische
Interpretation und ihre Grenzen werden in den Zustands- und Dynamikkapiteln
entwickelt.

\paragraph{Ausführungsvertrag.}
Diese Beziehungen empfangen realisierte metrische Größen und eine anwendbare
lokale Ingenieurnäherung; sie erzeugen Beschleunigungs-, Gleichgewichts- oder
Ruckgrößen. Sie setzen die angegebenen Geschwindigkeits- und
Differentiationsbedingungen voraus. Sie bewerten physikalische Folgen, können
jedoch weder betriebliche Zulässigkeit bestimmen noch eine Entscheidung
autorisieren.

\section{Track--World-Beziehungen}

In einer lokalen ebenen Realisierung wird ein gleisbezogener Punkt mit der
longitudinalen Koordinate \(s\) und dem seitlichen Versatz \(d\) durch
\begin{equation}
x_{\mathrm w}(s,d)
=
\gamma(s)+d\,\mathbf n(s)
\label{eq:track-to-world}
\end{equation}
in den World Space abgebildet.

\begin{definition}[Track-to-World-Abbildung]
\label{def:track-to-world}
Die durch \cref{eq:track-to-world} definierte Abbildung setzt gleisbezogene
Koordinaten zu einer World-Space-Realisierung in Beziehung.
\end{definition}

Das zugehörige inverse Problem wird ausgedrückt als
\begin{equation}
(s^\ast,d^\ast)
=
\operatorname*{arg\,min}_{s,d}
\left\|
x_{\mathrm w}-\bigl(\gamma(s)+d\,\mathbf n(s)\bigr)
\right\|.
\label{eq:world-to-track}
\end{equation}

\begin{definition}[World-to-Track-Abbildung]
\label{def:world-to-track}
Die World-to-Track-Abbildung bestimmt innerhalb eines anwendbaren Realization
Context eine intrinsische gleisbezogene Interpretation eines
World-Space-Punkts.
\end{definition}

Eine Transformation zwischen zwei World-Space-Koordinatenbeschreibungen wird
abstrakt geschrieben als
\begin{equation}
x_2=\mathcal C(x_1).
\label{eq:coordinate-transformation}
\end{equation}

Koordinatentransformation und World-to-Track-Abbildung sind verschiedene
Operationen: Erstere ändert die World-Space-Repräsentation, Letztere setzt den
World Space zu einer intrinsischen gleisbezogenen Domäne in Beziehung.

\paragraph{Ausführungsvertrag.}
Track-to-World empfängt ein realisiertes Alignment-Frame und gleisbezogene
Koordinaten; es erzeugt World-Koordinaten. World-to-Track empfängt einen
World-Punkt und dieselbe anwendbare Realisierung; es erzeugt eine mögliche
intrinsische Interpretation. Eindeutigkeit erfordert eine geeignete lokale
Domäne. Keine der beiden Operationen begründet, dass zwei Quellen dasselbe
Objekt betreffen.

\section{Repräsentation und ingenieurtechnische Realisierung}

Das Abtasten einer stetigen Funktion an den Positionen \(s_i\) wird
geschrieben als
\begin{equation}
\mathcal S(f)=\{f(s_i)\},
\label{eq:sampling}
\end{equation}
die Rekonstruktion aus abgetasteten Werten als
\begin{equation}
\mathcal I\bigl(\{f(s_i)\}\bigr)=\widetilde f(s).
\label{eq:interpolation}
\end{equation}

Ein ingenieurtechnischer Realisierungsoperator setzt einen beabsichtigten
Zustand zu einem physisch realisierten Zustand in Beziehung:
\begin{equation}
\mathcal R:
\mathrm{pose3}_{\mathrm{target}}
\longrightarrow
\mathrm{pose3}_{\mathrm{real}}.
\label{eq:realization-operator}
\end{equation}

Eine in späteren Realisierungskapiteln verwendete schematische Zerlegung lautet
\begin{equation}
\mathcal R
=
\mathcal M\circ\mathcal I\circ\mathcal S,
\label{eq:core-realization-decomposition}
\end{equation}
wobei Abtastung, Rekonstruktion und physische Antwort getrennte
Operatorverantwortungen bleiben.

Auf Krümmungsebene lautet die entsprechende Beziehung
\begin{equation}
\mathcal R\bigl(\kappa_{\mathrm{target}}\bigr)
=
\kappa_{\mathrm{real}}.
\label{eq:curvature-realization}
\end{equation}

\begin{definition}[Krümmungsrealisierung]
\label{def:curvature-realization}
Krümmungsrealisierung bezeichnet die Beziehung zwischen beabsichtigter und
physisch realisierter Krümmung unter einem anwendbaren ingenieurtechnischen
Realisierungsprozess.
\end{definition}

Das zugehörige inverse Problem sucht eine beabsichtigte Krümmung, die
\begin{equation}
\mathcal R\bigl(\kappa_{\mathrm{target}}\bigr)
\approx
\kappa_{\mathrm{desired}}
\label{eq:inverse-realization}
\end{equation}
erfüllt.

\paragraph{Ausführungsvertrag.}
Abtastung empfängt eine stetige Beschreibung und Abtastpositionen;
Interpolation empfängt abgetastete Werte und eine Interpolationsregel;
Realisierung empfängt einen Zielzustand und ein physisches Antwortmodell. Ihre
Ausgaben bewahren diese Abhängigkeiten. Im Testfall kann eine
Vermessungspolyline eine beobachtete Beschreibung realisierter Geometrie
stützen; weder Abtastung noch Interpolation macht die Beobachtung zur
physischen Realität oder verleiht ihr konstruktive Autorität.

\section{Flächenkrümmung}

Für eine an eine Fläche gebundene Kurve zerfällt die Gesamtkrümmung in einen
geodätischen und einen normalen Anteil:
\begin{equation}
\kappa^2=\kappa_g^2+\kappa_n^2.
\label{eq:curvature-decomposition}
\end{equation}

Die geodätische Krümmung wird intrinsisch ausgedrückt durch
\begin{equation}
\kappa_g
=
\left\|\nabla_s\mathbf t\right\|.
\label{eq:geodesic-curvature}
\end{equation}

Die anwendbare Fläche, Metrik, Verbindung und ingenieurtechnische
Interpretation werden in den Kapiteln über Ellipsoid und räumliche
Realisierung festgelegt.

\paragraph{Ausführungsvertrag.}
Diese Gleichungen empfangen eine Kurve, eine Fläche, eine Metrik und eine
Verbindung; sie erzeugen Krümmungsanteile in dieser Realisierung. Ohne diese
Eingaben sind die Symbole nicht ausführbar. Die Anteile können nicht
entscheiden, welches Flächenmodell für ein Projekt autoritativ ist.

\section{Optimierungsbeziehungen}

Ein schematisches realisierungsbewusstes Optimierungsproblem lautet
\begin{equation}
\min_u\;J\bigl(\mathcal R(x(u))\bigr).
\label{eq:general-optimization}
\end{equation}

Für nichtlineare Probleme der kleinsten Quadrate besitzt die Zielfunktion die
Form
\begin{equation}
\min_x\frac12\|r(x)\|^2,
\label{eq:least-squares}
\end{equation}
mit der Gauß--Newton-Näherung
\begin{equation}
H\approx J^T J.
\label{eq:gauss-newton}
\end{equation}

Diese Gleichungen spezifizieren lediglich verbreitete numerische Strukturen.
Ingenieurbeobachtungen, Residuen, Nebenbedingungen, Zulässigkeit und
Entscheidungen bewahren ihre vom Solver unabhängige Bedeutung.

\paragraph{Ausführungsvertrag.}
Optimierung empfängt eine deklarierte Zielfunktion, Variablen, Residuen,
Nebenbedingungen und jeden in das Modell aufgenommenen Realisierungsoperator;
sie erzeugt einen numerischen Kandidaten. Ihre Gültigkeit hängt von diesen
Entscheidungen und von Solver-Annahmen ab. Ein Minimum ist ein
Bewertungsergebnis, keine Ingenieurentscheidung.

\section{Kanonische Rolle}

Die Gleichungen bilden nun eine Abhängigkeitskette statt eines Formelkatalogs.
Sie können eine Entwurfstrajektorie rekonstruieren, sie unter metrischen
Qualifikationen realisieren, Beobachtungen interpretieren und qualifizierte
Differenzen berechnen. Die nächsten Kapitel stellen die Zustands- und
Operatormaschinerie bereit, die zur konsistenten Ausführung dieser Kette
erforderlich ist. Die Ingenieurbegriffe benennen anschließend die
Verantwortlichkeiten, die die Berechnung bewahren muss.
